\documentclass[11pt,a4paper]{article}
\usepackage[T1]{fontenc}
\usepackage{lmodern}
\usepackage[margin=26mm]{geometry}
\usepackage{amsmath,amssymb,amsthm,booktabs,microtype}
\usepackage[hidelinks]{hyperref}
\hypersetup{pdftitle={Pre-Branching Constraints for Max 2-SAT},
  pdfsubject={Sound learning, random-sign counting, and exact enumeration}}
\newtheorem{theorem}{Theorem}
\newtheorem{proposition}[theorem]{Proposition}
\newtheorem{lemma}[theorem]{Lemma}
\newtheorem{corollary}[theorem]{Corollary}
\theoremstyle{remark}
\newtheorem{remark}[theorem]{Remark}
\newcommand{\Ostar}{O^{*}}
\newcommand{\E}{\mathbb E}
\newcommand{\Prb}{\mathbb P}
\newcommand{\ind}[1]{\mathbf 1\{#1\}}
\emergencystretch=2em
\clubpenalty=10000
\widowpenalty=10000
\title{Pre-Branching Constraints for Max 2-SAT\\
\large Sound learning, simultaneous counting, and exact enumeration}
\author{}
\date{12 October 2026}
\begin{document}
\maketitle

\begin{abstract}
We give a common correctness invariant for constraints introduced before
branching: every comparison preserves the same optimal assignment, chosen
by a fixed secondary order. Single-variable comparisons yield guarded
cardinality constraints, and partial improvement certificates yield
ordinary clauses. Pair exchanges give a directly checkable source of
binary clauses and exact exponential savings when they contain a large
matching or long disjoint chains. For the random-sign model, we prove a
separate counting theorem: on any fixed interaction graph, with arbitrary
fixed initial unit clauses, the expected number of assignments satisfying
all single-variable comparisons is at most $2^{n/2}$. The proof uses a
read-two entropy inequality, not independence of the comparisons or any
degree bound. This is a candidate-count bound, not an enumeration-time
bound. We state the remaining Las Vegas enumeration target precisely.
The local derivations are proved here; literature novelty is not claimed.
\end{abstract}

\section{What is being constrained?}

Let $F(x)$ be the number of satisfied soft clauses of an unweighted
Max 2-SAT instance, including unit clauses. There are initially no genuine
hard feasibility constraints. All assignments belong to the original
optimization domain. We may also use the equivalent quadratic objective
\begin{equation}
 F(x)=c+\sum_u h_ux_u+\sum_{u<v}J_{uv}x_ux_v,
 \qquad x\in\{0,1\}^n.                         \label{eq:quadratic}
\end{equation}
Units affect $c,h$; binary clauses determine $J$. The deterministic rules
below also apply to rational coefficients of polynomial encoding length.

Fix a priority order on the variables and a preferred bit $a_u$ for each
variable. Rank complete assignments first by $F$, then lexicographically
by agreement with these preferred bits. Call the resulting strict total
order $\Phi$, and let $x^\star$ be its unique maximum. Thus $x^\star$ is
an optimum of the original problem. The priority and preferred bits are
fixed for all constraints that will be combined.

\begin{proposition}[Common representative]\label{prop:representative}
For any collection of maps $T:\{0,1\}^n\to\{0,1\}^n$, the constraints
\[
 \Phi(x)\ge \Phi(T(x))
\]
hold simultaneously at $x^\star$. Any sound consequences of these
constraints can therefore be added together without losing every optimum.
\end{proposition}
\begin{proof}
No assignment has rank greater than the maximum $x^\star$.
\end{proof}

The maps may change one variable, exchange two values, or replace a
specified pattern on a block. They may be selected adaptively or randomly.
We do not need to know $x^\star$ to construct the comparisons. Randomness
selects which valid tests to use; it does not guess that an unjustified
inequality holds. If the input has genuine hard constraints, a replacement
must additionally be feasible, or the comparison must be guarded by its
feasibility. That extension is not assumed here.

\paragraph{Compatibility with simplification.}
Strict domination for the soft objective is safe. A weak domination
fixing must respect the secondary order whenever equality is possible,
or have an additional compatibility proof. More generally, reductions
need a reconstruction invariant preserving $x^\star$ (or an alternative
joint proof). Existing path-dependent pruning constraints cannot simply
be mixed with a different convention for resolving ties.

For example, take $F=(u\lor v)$ and prefer zero at both variables. The
two single-variable comparisons below retain exactly $01$ and $10$.
Ordinary weak pure-literal domination can fix $u=1$ and then $v=1$
while preserving an optimum of the soft objective at each step; doing
so ignores the new constraints and eliminates both retained assignments.
This illustrates why separate preservation of \emph{some} optimum is
insufficient. A tie-aware implementation avoids this conflict.

\section{Single-variable constraints and short certificates}

Write
\[
 D_u(y)=F(u=1,y)-F(u=0,y)
       =h_u+\sum_{v\ne u}J_{uv}y_v.
\]
For readability, first take $a_u=1$. The full pre-branching constraint is
\begin{equation}
 C_u:\quad
 (u=1\Rightarrow D_u(y)\ge0)
 \ \land\ (u=0\Rightarrow D_u(y)\le-1).
                                                        \label{eq:guard}
\end{equation}
For $a_u=0$, use $D_u\ge1$ on the $u=1$ side and $D_u\le0$ on the
$u=0$ side. These thresholds use integer scores; strict inequalities
give the corresponding rational-score formulation. The guards involving
$u$ remain present until its value is fixed.

In the occurrence notation, let
\[
 k=d_1(u)-d_1(\neg u),\quad r=d_2(u),\quad q=d_2(\neg u),
 \quad G=r+k-q.
\]
Let $A$ count true partner literals of $u$ and $B$ true partner literals
of $\neg u$, counting occurrences. Then
\begin{equation}
 D_u=G-A+B=L-(q-k),\qquad L=(r-A)+B.             \label{eq:gap}
\end{equation}
Thus, with ties favoring $u=1$, the two guarded inequalities are
\begin{align}
 u=1&\Rightarrow L\ge q-k,\nonumber\\
 u=0&\Rightarrow L\le q-k-1.                    \label{eq:two}
\end{align}
The second is the usual pruning constraint. Its one-neighborhood
consequence is
\begin{equation}
 u=0\Rightarrow A\ge G+1.                      \label{eq:one}
\end{equation}
This uses the user's occurrence gap: unit and binary occurrences count
equally. It does not change that definition to the differently weighted
quantity used to select a branching orientation in some earlier notes.

The reverse upper-bound form, obtained by complementing every pruning
literal, is
\[
 u=1\Rightarrow A+(q-B)\le r+k.
\]
The two full constraints allocate ties to exactly one value. Requiring
strict improvement on both sides is unsound. Dropping a nonnegative
term gives a valid weaker consequence; independently splitting an
allowance into smaller budgets does not.

\begin{proposition}[Partial certificates become clauses]\label{prop:certificate}
Let $\alpha$ be a partial assignment not fixing $u$. If every completion
of $\alpha$ has $D_u>0$, or has $D_u\ge0$ and $a_u=1$, then the ordinary
clause $\neg\alpha\lor u$ is sound under Proposition~\ref{prop:representative}.
The analogous negative certificate gives $\neg\alpha\lor\neg u$.
\end{proposition}
\begin{proof}
Any assignment violating the clause improves in $\Phi$ after flipping $u$.
\end{proof}

The certificate is efficiently checkable. Substitute $\alpha$ into $D_u$,
leaving $b+\sum_{v\in R}c_vy_v$. Its exact unconstrained range is
\begin{equation}
 \left[b+\sum_{v\in R}\min(0,c_v),\quad
       b+\sum_{v\in R}\max(0,c_v)\right].        \label{eq:range}
\end{equation}
This test may ignore older hard pruning constraints and is still sound.
Using them to obtain a stronger certificate requires a correct joint test.

For the simple-clause case, if $p=q-k\ge1$, any $p$ distinct pruning
literals of $L$ that are all true certify $u=1$ (ties favoring one).
They give a clause of width $p+1$ including $u$. When $r+k\ge0$, for reverse
strict improvement, $r+k+1$ true literals in the complemented list certify
$u=0$, giving width $r+k+2$. Such witnesses must actually exist; the
threshold can exceed the list length. Nonpositive certificate thresholds
are handled directly by the range test and can give a unit clause.
Multiplicities or opposite literals
on one variable require the occurrence-aware range test above. A full
cardinality constraint can represent exponentially many such clauses.

\begin{proposition}[One bit per full comparison]\label{prop:half}
For every fixed formula and variable $u$, exactly $2^{n-1}$ complete
assignments satisfy $C_u$.
\end{proposition}
\begin{proof}
For each assignment to the other variables, exactly one value of $u$
is preferred, with the specified tie convention.
\end{proof}

Consequently, a larger literal gap does not improve the unconditional
one-bit saving of the \emph{full guarded comparison}. Its importance is
that the preferred value can be certified after fewer neighboring values
are exposed, or by a shorter learned clause. Joint information and
early detectability are different quantities.

\section{A deterministic counting baseline: independent centers}

\begin{proposition}\label{prop:independent}
Let $T$ contain $s$ variables with $J_{uv}=0$ for all distinct $u,v\in T$.
Exactly $2^{n-s}$ assignments satisfy $\bigwedge_{u\in T}C_u$.
They can be generated and scored in $\Ostar(2^{n-s})$ time and
polynomial space.
\end{proposition}
\begin{proof}
Assign the outside variables. Each center's comparison now determines
its unique preferred value independently of the other centers. Enumerate
the $2^{n-s}$ outside assignments and compute these responses.
\end{proof}

The external neighborhoods may coincide completely and have linear size.
Repeated information can fix different center variables. This is the
familiar independent-set mechanism; balanced bipartite interaction
graphs give $\Ostar(2^{n/2})$ time. It is a useful baseline rather than
a new general algorithm. Identification of $T$ is part of the algorithm:
the proposition assumes it is supplied or efficiently found.

\section{A counting theorem for arbitrary graphs with random signs}
\label{sec:random}

Fix any simple graph $G=(V,E)$, any initial unit clauses, and the secondary
order, all before drawing the binary-clause signs. For each edge, choose
one of its four literal-sign pairs uniformly and independently. Let
\[
 Z(F)=\left|\{x:\ C_u(x)\text{ holds for every }u\}\right|.
\]
There is no restriction on minimum degree, maximum degree, regularity,
or the number of fixed initial units.

\begin{theorem}[Read-two count]\label{thm:random}
For each vertex $u$, let $p_u$ be the probability, over the clause signs,
that its preferred response is one, for any fixed assignment to the
other variables. This probability is independent of that assignment.
Then
\begin{equation}
 \boxed{\E_F Z(F)
 \le \prod_u\left(\sqrt{p_u}+\sqrt{1-p_u}\right)
 \le 2^{n/2}.}                                  \label{eq:randomcount}
\end{equation}
In particular, for every fixed $0<\varepsilon<1/2$,
\begin{equation}
 \Prb_F\left[Z(F)>2^{(1/2+\varepsilon)n}\right]
 \le 2^{-\varepsilon n}.                        \label{eq:randomtail}
\end{equation}
\end{theorem}

\begin{proof}
Fix a full assignment $x$. Write $b_u=d_1(u)-d_1(\neg u)$ and
$d_u=\deg_G(u)$. The distribution of $D_u(x_{-u})$ is
\begin{equation}
 b_u+\sum_{j=1}^{d_u}W_j,\qquad
 \Prb[W_j=-1]=\Prb[W_j=1]=\tfrac14,\quad
 \Prb[W_j=0]=\tfrac12.                          \label{eq:lazy}
\end{equation}
Indeed, a clause contributes zero when its partner literal is true and
otherwise contributes the sign of its literal on $u$. Independent uniform
incidence signs make this distribution independent of $x_{-u}$.
Thus $\Prb_F[C_u(x)]=p_u$ if $x_u=1$, and $1-p_u$ if $x_u=0$.

The events $C_u(x)$ are not independent. However, each independent
four-valued edge label is used only by the comparisons at its two
endpoints. The read-two form of the generalized H\"older/Shearer
inequality~\cite{Finner,ReadK} gives
\begin{equation}
 \Prb_F\left[\bigwedge_u C_u(x)\right]
 \le \prod_u\Prb_F[C_u(x)]^{1/2}.              \label{eq:readtwo}
\end{equation}
For completeness, in this finite uniform setting it follows directly
from entropy. Let $\mathcal E$ be the set of edge-label assignments
satisfying all the events, and $\mathcal A_u$ the allowed assignments
to the $d_u$ incident edge labels. If $\mathcal E$ is nonempty, apply
Shearer's entropy inequality to the uniform distribution on $\mathcal E$.
Every edge coordinate is covered twice, so
\[
 2\log_2|\mathcal E|
 \le\sum_u\log_2|\mathcal A_u|
 =4|E|+\sum_u\log_2\Prb_F[C_u(x)].
\]
Subtracting $4|E|$ proves~\eqref{eq:readtwo}; an empty event is trivial.
Summing over $x$ factorizes the right-hand side:
\[
 \E_F Z(F)
 \le\sum_x\prod_u
 \begin{cases}\sqrt{p_u},&x_u=1,\\\sqrt{1-p_u},&x_u=0\end{cases}
 =\prod_u(\sqrt{p_u}+\sqrt{1-p_u}).
\]
Each factor is at most $\sqrt2$. Markov's inequality yields
\eqref{eq:randomtail}.
\end{proof}

The inequality of Shearer used above follows from the entropy chain rule:
expand the entropy of every coordinate subset in the same global order;
conditioning on only the earlier coordinates of that subset gives at
least the corresponding conditional entropy given all earlier coordinates;
then sum, using the fact that each coordinate occurs twice.

The marginal can be computed exactly. With ties favoring one,
\[
 p_u=4^{-d_u}\sum_{\substack{-d_u\le s\le d_u\\b_u+s\ge0}}
                  \binom{2d_u}{d_u+s}.
\]
Replace $\ge0$ by $>0$ when ties favor zero. This includes isolated
variables. Their comparison fixes the preferred optimal value, so they
contribute a factor one, not two. Without a tie convention, the empty
formula would have $2^n$ weak local optima and the theorem would fail.

\paragraph{What the theorem establishes.}
For instance, take $\varepsilon=1/4$. With probability at least
$1-2^{-n/4}$ the entire constraint system has at most
$2^{3n/4}\simeq1.682^n$ solutions. The expectation has the stronger
base $\sqrt2$. These are input-distribution statements on \emph{counts}.
They do not yet give these running-time bounds. The signs must be
independent at clause incidences; randomly complementing the names of
variables in a fixed adversarial formula does not create that model.
The tie convention must also be fixed independently of the signs for
the stated marginal factorization.

If comparisons are imposed only for a fixed set $T$, the same proof,
with unrestricted local events at the other vertices, gives
\[
 \E_F\left|\{x:\bigwedge_{u\in T}C_u(x)\}\right|
 \le 2^{n-|T|}\prod_{u\in T}(\sqrt{p_u}+\sqrt{1-p_u})
 \le 2^{n-|T|/2}.
\]
The proof requires neither spectral assumptions nor a bound on overlap.
This is the principal counting result of this note. Local maxima are a
classical object in exact Max 2-SAT algorithms~\cite{Hastings}; a focused
literature comparison is still needed to establish novelty of this
particular formulation and bound.

\section{Pair exchanges give short pre-branching clauses}
\label{sec:exchange}

Single-variable comparisons need not be the only permissible constraints.
Compare the states $(x_u,x_v)=(1,0)$ and $(0,1)$. Their score difference is
\begin{equation}
 E_{uv}(y)=h_u-h_v+
       \sum_{w\notin\{u,v\}}(J_{uw}-J_{vw})y_w. \label{eq:exchange}
\end{equation}
The internal edge $uv$ cancels. Its exact minimum over all outside
assignments is
\begin{equation}
 \lambda_{uv}=h_u-h_v+
       \sum_{w\notin\{u,v\}}\min(0,J_{uw}-J_{vw}).
                                                        \label{eq:exmin}
\end{equation}

\begin{proposition}[Unconditional exchange clause]\label{prop:exchange}
If $\lambda_{uv}>0$, or if $\lambda_{uv}=0$ and the secondary order
prefers state $(1,0)$ to $(0,1)$, then the binary clause
\begin{equation}
 \boxed{x_u\lor\neg x_v}\qquad(x_v\Rightarrow x_u) \label{eq:implication}
\end{equation}
can be added before branching. Any collection passing these tests with
the common secondary order preserves $x^\star$.
\end{proposition}
\begin{proof}
Every violating assignment has $(x_u,x_v)=(0,1)$. Exchanging these values
increases $\Phi$, by~\eqref{eq:exchange}--\eqref{eq:exmin}.
\end{proof}

Identical interaction rows outside the pair reduce the test to comparing
$h_u$ and $h_v$. More generally the field advantage need only compensate
for the total unfavorable row discrepancy:
\[
 h_u-h_v\ge
 \sum_{w\notin\{u,v\}}\max(0,J_{vw}-J_{uw}),
\]
with the stated tie condition. Hence exact twins are unnecessary.
Both variables can still be individually undominated. All pair tests
are computable in $O(n^3)$ arithmetic operations using a dense matrix.
Variables may be complemented once before applying the test.

After exposing a partial assignment $\alpha$, the same minimum test on
the residual expression gives the guarded clause
$\neg\alpha\lor x_u\lor\neg x_v$. This is one concrete way that branching
on discrepancy coordinates creates new short constraints.

\begin{theorem}[Matching and chain savings]\label{thm:chains}
Suppose the valid unconditional exchange clauses contain $r$ clauses
on pairwise disjoint variable pairs. They admit exactly
\begin{equation}
 3^r2^{n-2r}=2^n(3/4)^r                         \label{eq:matching}
\end{equation}
assignments. Enumerating these assignments solves the original soft
instance in this time up to polynomial factors, with polynomial space.

More generally, suppose selected clauses form disjoint chains on
$s_1,\ldots,s_q$ variables, requiring ones to precede zeros in each chain.
Then the corresponding bound is
\begin{equation}
 2^{n-\sum_i s_i}\prod_i(s_i+1)=2^{n-\Gamma},
 \qquad \Gamma=\sum_i\bigl(s_i-\log_2(s_i+1)\bigr). \label{eq:chains}
\end{equation}
This too is an exact polynomial-space running-time bound.
\end{theorem}
\begin{proof}
A pair has three allowed states. A chain of length $s_i$ has $s_i+1$
allowed cut positions. Enumerate these choices and the remaining free
bits, and score every generated assignment. There is no dead-prefix
overhead: every generated prefix extends to a state of this relaxation.
Its domain contains $x^\star$, even if additional valid constraints are
ignored during generation.
\end{proof}

A linear matching gives a fixed exponential saving independent of
clause density. A maximum matching in the certified-pair graph can be
found in polynomial time; even a greedy maximal matching loses at most
a factor two in matching size. Thus this is a checkable class, not a
condition involving unknown numbers of satisfying assignments.
Longer chains use overlap constructively: a chain on all variables
reduces the candidate set to $n+1$ assignments.

\paragraph{A dense redundant example that is resolved.}
For even $n$, put $(x_u\lor x_v)$ and $(\neg x_u\lor\neg x_v)$ on
every pair. If $k$ bits are one, the score is
\[
 \binom n2+k(n-k).
\]
All $\binom n{n/2}$ balanced assignments are strict single-flip local
optima, so every choice of single-variable tie preferences retains them.
This standard example is discussed in~\cite{Hastings}. But its exchange
differences vanish identically. Choose a priority order preferring ones;
the adjacent exchange clauses impose one chain, leaving $n+1$ assignments,
and only one survives the single-variable comparisons as well.
The redundant family therefore suggests an additional learning rule
rather than an impossibility for the algorithmic program.

The matching rule overlaps with the existing near-twin and group-count
notes~\cite{Near,Groups}. The contribution here is its explicit encoding
as a compatible pre-branching clause, its discrepancy-versus-field test,
and its integration with counting. It does not subsume the group-count
solver, whose cuts allow other forms of arbitrary internal structure.

\section{General block certificates and their counting effect}

Fix a block $S$, an excluded pattern $a\in\{0,1\}^S$, and a replacement
$b\in\{0,1\}^S$. Quadraticity implies
\begin{align}
 F(b,y)-F(a,y)&=\kappa+\sum_{w\notin S}c_wy_w,\nonumber\\
 \kappa&=\sum_{u\in S}h_u(b_u-a_u)
       +\sum_{\{u,v\}\subseteq S}J_{uv}(b_ub_v-a_ua_v),\nonumber\\
 c_w&=\sum_{u\in S}J_{uw}(b_u-a_u).              \label{eq:block}
\end{align}
Thus $\kappa+\sum_w\min(0,c_w)$ is an exact, polynomial-time test
for unconditional domination by this \emph{specified} replacement.
Strict positivity, or zero with a favorable secondary comparison,
allows the ordinary nogood clause
\[
 \bigvee_{u\in S}(x_u\ne a_u).
\]
The comparison of the two block patterns in the secondary order is
independent of the outside assignment. Searching over all blocks or all
replacements is not claimed to be polynomial.

For disjoint supports of widths $w_1,\ldots,w_t$, one forbidden pattern
per support leaves exactly
\[
 2^n\prod_{i=1}^t(1-2^{-w_i})
\]
assignments and gives an enumeration algorithm of that cost. Bounded
width and linearly many such clauses suffice. Growing width alone gives
no comparable constant saving from this product; overlap can help, as
the chain example shows, but must be analyzed jointly.

\section{Enumeration and Las Vegas targets}

For a supplied exposed set $K$, suppose polynomial-time propagation of
compatible constraints fixes at least $\delta n$ variables outside $K$
on all but a $2^{-\gamma n}$ fraction of assignments to $K$ (infeasible
assignments can be discarded). Assume $|K|\le(1-\delta)n$.
Enumerating $K$ and then the remaining free bits costs
\begin{equation}
 \Ostar\left(2^{n-\delta n}+2^{n-\gamma n}\right). \label{eq:exposure}
\end{equation}
This conditional theorem pays for all exposed assignments. Its missing
ingredient in a new formula class is availability of such a $K$ and the
stated propagation bound, not the final summation.

Alternatively, randomize an exact search strategy. If an independent
attempt finishes within $B$ operations with probability at least $p>0$,
abort unfinished attempts at $B$ and restart. Returning only after a
complete exact search is a Las Vegas algorithm with expected time
$O(B/p)$. In particular,
\[
 B=\Ostar(2^{(1-\delta)n}),\quad p\ge2^{-\eta n},\quad \eta<\delta
 \quad\Longrightarrow\quad
 \E T=\Ostar(2^{(1-\delta+\eta)n}).
\]
Polynomial space is retained if each attempt has polynomial space and
discarded attempt data are not accumulated without a bound. Keep the
secondary convention fixed when retaining constraints across attempts.

\paragraph{What a direct random-order analysis must bound.}
For a fixed constraint system and a monotone rejection test, assign
variables in a uniformly random permutation. Let $q_t$ be the probability
that a uniform assignment to its first $t$ variables passes the test.
The expected number of passing depth-$t$ prefixes is $2^tq_t$, and the
expected traversal cost, up to polynomial factors, is bounded by
\[
 1+\sum_{t=0}^n2^tq_t.
\]
At depth $n$, if the test evaluates the whole constraint system,
$2^nq_n=Z(F)$. A candidate-count theorem controls this final term.
A direct search proof also controls earlier terms; a restart proof may
instead establish a sufficiently frequent set of fast attempts.

For the random-sign class, Theorem~\ref{thm:random} settles a strong
counting target for the pre-branching constraints. Obtaining a Las Vegas
enumerator with a comparable base remains open here. The earlier
early-enumeration note~\cite{Early} proves a different, weaker-base time
bound under almost regularity by enumerating a tractable summed
relaxation. That theorem is not automatically upgraded to $(\sqrt2)^n$
by the count proved above.

\section{Checks and the next mathematical questions}

The companion \texttt{explore\_prebranch\_constraints.py} uses integer
scores and exact rational probabilities. Its output is
\texttt{prebranch\_constraint\_checks.json}. All checks passed:
\begin{itemize}
\item 100 formulas on 4--7 variables, 37,600 exact flip identities,
6,339 one-neighborhood implications, and preservation of the canonical
optimum under simultaneous single-variable and exchange constraints.
\item 12,204 exchange implications; 100 independent-center counts and
100 matching counts; 5,364 block-minimum checks and 1,547 block nogoods.
\item Exhaustive enumeration of 12,497 signings across eight graph/unit/tie
models, including all $4^6$ signings of $K_4$ in three settings. Every
fixed-assignment read-two inequality was checked using rational arithmetic.
\item Complete-cut examples on $n=6,8,10,12$ variables: single-variable
counts $20,70,252,924$ respectively; chain counts $7,9,11,13$;
their conjunction has one assignment in each case.
\end{itemize}
These finite checks are sanity checks for the stated proofs, not
asymptotic evidence or a substitute for them.

The existing exploratory command \texttt{python3 explore\_two\_sided.py combine}
also provides a concrete compatibility warning. On its 300 test instances,
mixing constraints registered with a fixed root orientation and
path-dependent strict branching constraints lost the optimum in 17 cases.
The branching-only baseline and the variant whose branching followed the
registered orientation lost none in that run. This experiment is not an
implementation of the full common-order system above; it confirms that
an unproved mixture of conventions is a real correctness risk.

\newpage
The companion script also checks
the auxiliary identity
\[
 \E_a Z(F,a)=\sum_{x\text{ a weak single-flip local optimum}}2^{-z(x)},
\]
where $a$ is a vector of independent uniform tie preferences and $z(x)$
counts zero-gain single flips at $x$. Each tied coordinate imposes one
independent preferred-bit condition. This identity shows exactly what
random tie choices can remove, and why strict local optima need other
comparisons.

The next questions are concrete:
\begin{enumerate}
\item For random-sign inputs, can an exposure strategy or a counted
relaxation enumerate the pre-branching comparison solutions without
visiting almost all dead prefixes?
\item When exchange tests initially fail, can a small exposure set make
a linear matching of tests succeed on enough branches, with the cost
of the other branches explicitly included?
\item Can persistent redundancy of the comparisons be converted into
chains, collective block certificates, or the group-count structures of
the earlier note, with a uniform bound on the total search cost?
\end{enumerate}
The local soundness rules, random-sign counting theorem, and matching/chain
algorithms above are proved. A general availability theorem joining
these alternatives is not claimed.

\begin{thebibliography}{9}
\bibitem{Finner} H.~Finner. A generalization of H\"older's inequality and
some probability inequalities. \emph{Annals of Probability} 20(4),
1893--1901, 1992. \url{https://doi.org/10.1214/aop/1176989534}.
\bibitem{ReadK} D.~Gavinsky, S.~Lovett, M.~Saks, S.~Srinivasan.
A tail bound for read-$k$ families of functions. 2012 preprint.
\url{https://arxiv.org/abs/1205.1478}.
\bibitem{Hastings} M.~B.~Hastings. Local maxima and improved exact algorithm
for MAX-2-SAT. 2016. \url{https://arxiv.org/abs/1610.07100}.
\bibitem{Near} \emph{Near-Twin Neighborhoods as a Structural Exit}.
Working note, 8 October 2026. \path{max2sat_near_twin_pairs.tex}.
\bibitem{Groups} \emph{Useful Constraints or Compression by Group Counts}.
Working note, 10 October 2026.
\path{max2sat_redundancy_and_group_counts.tex}.
\bibitem{Early} \emph{Max 2-SAT by Early Constraint-Guided Enumeration}.
Working note, 8 October 2026. \path{max2sat_early_enumeration.tex}.
\end{thebibliography}
\end{document}
